\section{Threats to Validity and Conclusion}
\label{sec:conclusion}

\paragraph{Threats to validity.} The search may miss grey literature and vendor work, and
arXiv records are not peer reviewed; snowballing and labelling mitigate but do not remove
this. AUF scores are first-rater codes based on default specifications until the second
rater's agreement is reported, and optional extensions or deployed profiles can change
judgements. The capability list was derived top-down and is tested for saturation against
the coded corpus. The case study uses one authored plant whose tasks, faults and
validators share a specification; its OPC~UA/AAS comparison measures emitted profiles from
that specification. No physical trial or complete model evaluation is reported; the
surrogate tokenizer does not measure billing tokens.

\paragraph{Conclusion.} Industrial AI-agent reliability depends less on the choice of
language model or semantic syntax than on whether the agent has access to explicit,
current, machine-interpretable engineering evidence covering identity, quantities,
relations, behaviour, constraints and provenance. Existing standards provide many of
these capabilities, but they are fragmented across technologies and rarely exposed as a
coherent runtime evidence contract. The case study shows the consequences: typing alone
resolves targets but not admissibility, equally informed validators decide identically,
and matched \opcua{} and AAS profiles are complementary yet jointly incomplete. Making
industrial standards agent-ready, chiefly by publishing behaviour and constraints with
provenance behind query interfaces, is the most consequential next step.
